% =============================================================================
% preambulo.tex - Configuracoes Tipograficas e Pacotes (Conforme ABNT NBR 14724)
% =============================================================================

\usepackage{cmap}
\usepackage[T1]{fontenc}
\usepackage[utf8]{inputenc}

% Margens estritas da ABNT NBR 14724 (formato anverso / digital)
% Superior: 3 cm (cabecalho a 2 cm com headsep=10mm)
% Esquerda: 3 cm
% Inferior: 2 cm
% Direita:  2 cm (numero a 2 cm da borda direita)
\usepackage[
    a4paper,
    top=30mm,
    bottom=20mm,
    left=30mm,
    right=20mm,
    headheight=15pt,
    headsep=10mm,
    footskip=10mm
]{geometry}

% Tipografia padronizada Times New Roman (NBR 14724 item 5.1)
\usepackage{newtxtext}
\usepackage{newtxmath}
\usepackage[scaled=0.88]{inconsolata}

\usepackage{lastpage}
\usepackage{indentfirst}
\usepackage{microtype}
\usepackage{graphicx}
\usepackage{subcaption}
\usepackage{booktabs}
\usepackage{multirow}
\usepackage{makecell}

% Matematica avancada
\usepackage{mathtools}
\usepackage{bm}

% Unidades do SI (configuradas para o padrao brasileiro / ABNT)
\usepackage{siunitx}
\sisetup{
    output-decimal-marker={,},
    group-separator={.},
    inter-unit-product=\ensuremath{{\cdot}}
}
\DeclareSIUnit\calorie{cal}
\DeclareSIUnit\molar{mol}
\DeclareSIUnit\litre{L}

% Graficos vetoriais de alta precisao com TikZ e PGFPlots
\usepackage{tikz}
\usepackage{pgfplots}
\pgfplotsset{compat=1.18}
\usetikzlibrary{shapes.geometric,arrows.meta,positioning,calc,fit,patterns}

% Cores para codigos e diagramas
\usepackage[dvipsnames]{xcolor}

% Formatacao profissional de codigo-fonte (listings)
\usepackage{listings}
\renewcommand{\lstlistingname}{Código}
\renewcommand{\lstlistlistingname}{Lista de Códigos}

\lstdefinestyle{pythonstyle}{
    language=Python,
    basicstyle=\ttfamily\footnotesize,
    columns=flexible,
    keepspaces=true,
    keywordstyle=\color{NavyBlue}\bfseries,
    commentstyle=\color{ForestGreen}\itshape,
    stringstyle=\color{BrickRed},
    numberstyle=\tiny\color{gray},
    numbers=left,
    numbersep=8pt,
    frame=single,
    framerule=0.5pt,
    rulecolor=\color{gray!30},
    backgroundcolor=\color{gray!4},
    breaklines=true,
    breakatwhitespace=true,
    showstringspaces=false,
    captionpos=b,
    tabsize=4,
    xleftmargin=12pt,
    framexleftmargin=12pt
}
\lstset{style=pythonstyle}

% Citacoes no padrao ABNT (carrega hyperref internamente)
\usepackage[alf]{abntex2cite}

% Configuracao de links e metadados no PDF
\makeatletter
\hypersetup{
    pdftitle={\@title},
    pdfauthor={Gabriel Marques de Andrade; Elilton Rodrigues Edwards},
    pdfsubject={Relatório Técnico-Científico PROIC/FAPESB/CNPq 34/2026},
    pdfcreator={LaTeX with abnTeX2},
    colorlinks=true,
    linkcolor=NavyBlue,
    citecolor=NavyBlue,
    urlcolor=NavyBlue
}
\makeatother

% Cleveref para referencias cruzadas em portugues
\usepackage[nameinlink,brazilian]{cleveref}
\crefname{equation}{Eq.}{Eqs.}
\crefname{figure}{Figura}{Figuras}
\crefname{table}{Tabela}{Tabelas}
\crefname{chapter}{Capítulo}{Capítulos}
\crefname{section}{Seção}{Seções}
\crefname{subsection}{Subseção}{Subseções}
\crefname{lstlisting}{Código}{Códigos}

% Padronizacao tipografica dos titulos de capitulos e secoes (NBR 6024 e NBR 14724)
\renewcommand{\ABNTEXchapterfont}{\normalfont\bfseries}
\renewcommand{\ABNTEXchapterfontsize}{\Large}
\renewcommand{\ABNTEXsectionfont}{\normalfont\bfseries}
\renewcommand{\ABNTEXsectionfontsize}{\large}
\renewcommand{\ABNTEXsubsectionfont}{\normalfont\bfseries}
\renewcommand{\ABNTEXsubsectionfontsize}{\normalsize}
\renewcommand{\ABNTEXsubsubsectionfont}{\normalfont\itshape}
\renewcommand{\ABNTEXsubsubsectionfontsize}{\normalsize}
\renewcommand{\ABNTEXsubsubsubsectionfont}{\normalfont}
\renewcommand{\ABNTEXsubsubsubsectionfontsize}{\normalsize}

% Parametrizacao de paragrafo e espacamento ABNT NBR 14724
\setlength{\parindent}{1.25cm}
\setlength{\parskip}{0pt}

% Formatacao de legendas e fontes (ABNT NBR 14724 itens 5.8 e 5.9: fonte menor 10pt)
\captiondelim{ -- }
\captionnamefont{\footnotesize\bfseries}
\captiontitlefont{\footnotesize}

% Estilo de pagina textual ABNT: numero no canto superior direito, sem cabecalho
\makepagestyle{abnttextual}
\makeevenhead{abnttextual}{}{}{\footnotesize\thepage}
\makeoddhead{abnttextual}{}{}{\footnotesize\thepage}
\makeevenfoot{abnttextual}{}{}{}
\makeoddfoot{abnttextual}{}{}{}

% Redefinicao da Capa no padrao institucional UESC / ABNT NBR 14724
\renewcommand{\imprimircapa}{%
  \begin{capa}%
    \center
    {\ABNTEXchapterfont\bfseries\normalsize\MakeTextUppercase{\imprimirinstituicao}\par}
    \vspace*{3.0cm}
    {\ABNTEXchapterfont\bfseries\normalsize\MakeTextUppercase{\imprimirautor}\par}
    \vspace*{\fill}
    {\ABNTEXchapterfont\bfseries\large\MakeTextUppercase{\imprimirtitulo}\par}
    \vspace*{\fill}
    {\ABNTEXchapterfont\bfseries\normalsize\MakeTextUppercase{\imprimirlocal}\par}
    {\ABNTEXchapterfont\bfseries\normalsize\imprimirdata\par}
    \vspace*{1cm}
  \end{capa}
}

% Redefinicao da Folha de Rosto no padrao ABNT NBR 14724
\renewcommand{\imprimirfolhaderosto}{%
  \clearpage
  \begin{center}
    {\ABNTEXchapterfont\bfseries\normalsize\MakeTextUppercase{\imprimirautor}\par}
    \vspace*{\fill}
    {\ABNTEXchapterfont\bfseries\large\MakeTextUppercase{\imprimirtitulo}\par}
    \vspace*{\fill}
    \hspace{.45\textwidth}
    \begin{minipage}{.50\textwidth}
      \SingleSpacing
      \footnotesize
      \imprimirpreambulo
      \par\vspace{1.5em}
      {\textbf{Orientador:} \imprimirorientador\par}
    \end{minipage}%
    \vspace*{\fill}
    {\ABNTEXchapterfont\bfseries\normalsize\MakeTextUppercase{\imprimirlocal}\par}
    {\ABNTEXchapterfont\bfseries\normalsize\imprimirdata\par}
    \vspace*{1cm}
  \end{center}
  \clearpage
}

% Comandos matematicos personalizados
\newcommand{\ddt}{\frac{\mathrm{d}}{\mathrm{d}t}}
\newcommand{\deriv}[2]{\frac{\mathrm{d}#1}{\mathrm{d}#2}}
\newcommand{\CA}{C_A}
\newcommand{\Tss}{T_{\mathrm{ss}}}
\newcommand{\qj}{q_j}
\newcommand{\Tj}{T_j}
\newcommand{\Tjf}{T_{jf}}
\newcommand{\Caf}{C_{Af}}
\newcommand{\rA}{r_A}
\newcommand{\UAnom}{UA_{\mathrm{nom}}}
\newcommand{\UAeff}{UA_{\mathrm{eff}}}
